\section{Methods and protocol}
\label{sec:methods}
Protocol version 0.1 is dated October 21, 2026 and has not been registered.
The full operational protocol is maintained in \texttt{protocol.md}.
Eligible populations are adults with diagnosed ADHD; mixed samples require separately extractable adult-ADHD findings.
Interventions are structured CBT programs, with cognitive restructuring, organization, planning, therapist support and delivery format extracted separately.
Comparators are classified as passive, active psychological, or medication-only for adjunctive treatment.
Primary clinical outcomes are ADHD symptoms and daily functioning, separated by rater and time point.

The meta-review includes systematic reviews with explicit search and selection methods.
The supplementary original-study stream requires pretreatment executive-function measurement and a later outcome association or treatment interaction.
Single-arm associations are prognostic evidence, not causal treatment-effect moderation.
Post-treatment executive-function scores alone do not enter the baseline moderator synthesis.
We will record whether measures are performance-based or self-reported, as well as diagnosis, medication, comorbidity, outcome overlap, missingness and prespecification.

Reviews will be appraised using AMSTAR 2, with domain judgments retained \citep{Shea2017}.
The original-study risk-of-bias instrument will be fixed in a protocol amendment before extraction.
Review overlap will be mapped at the primary-study level; overlapping meta-analyses will not be pooled as independent estimates.
Quantitative moderation synthesis remains conditional on compatible constructs, comparisons and interaction estimands, with a statistical analysis amendment required before fitting.
